\appendix
\subsection{Derivation of the Optimal Next Proposal}
\label[appendix]{sec:appendix:lagrangian}

Using the \ac{KL} divergence and entropy, given by
\begin{align}
    D_{\mathrm{KL}}(\pdf_{j+1} \| \pdf_{j})
    &= \int_{\IR^{\dimVarCem}} \pdf_{j+1} (\vVarCem) \log \pdf_{j+1} (\vVarCem) \dd \vVarCem
    \\
    &\qquad
    - \int_{\IR^{\dimVarCem}} \pdf_{j+1} (\vVarCem) \log \pdf_{j} (\vVarCem) \dd \vVarCem \enspace, \\
    H(\pdf_{j+1})
    &= -\int_{\IR^{\dimVarCem}} \pdf_{j+1} (\vVarCem) \log \pdf_{j+1} (\vVarCem) \dd \vVarCem \enspace,
\end{align}
the expanded Lagrangian~\cref{eq:tr:lagrangian} reads
\begin{align}
    &\mathcal{L}(\pdf_{j+1},\eta,\alpha, \nu)
    = \int_{\IR^{\dimVarCem}} \pdf_{j+1}(\vVarCem) J(\vVarCem) \, \dd \vVarCem \\
    &\qquad + (\eta + \alpha) \int_{\IR^{\dimVarCem}} \pdf_{j+1}(\vVarCem) \log \pdf_{j+1}(\vVarCem) \dd \vVarCem \\
    &\qquad - \eta \int_{\IR^{\dimVarCem}} \pdf_{j+1}(\vVarCem) \log \pdf_{j}(\vVarCem) \dd \vVarCem \\
    &\qquad - \eta \epsilon + \alpha H_{\min} 
    % \\
    % & \qquad 
    + \nu\left(\int_{\IR^{\dimVarCem}} \pdf_{j+1}(\vVarCem) \dd  \vVarCem - 1\right) \enspace.
\end{align}
Taking the functional derivative point-wise w.r.t. $\pdf_{j+1}$ yields
\begin{align}
    \label{eq:lagrange_func_derivative}
    \frac{\delta \mathcal{L}}{\delta \pdf_{j+1}}
    & = J(\vVarCem) + (\eta + \alpha) \left( \log \pdf_{j+1}(\vVarCem) + 1 \right) 
    \\
    & \qquad
    - \eta \log \pdf_{j}(\vVarCem) + \nu \enspace.
\end{align}
Setting this derivative to zero and solving for $\log \pdf_{j+1}$ gives
\begin{align}
    \label{eq:log_opt_proposal}
    \log \pdf_{j+1}^{\ast}(\vVarCem)
    = -\frac{J(\vVarCem)}{\eta + \alpha}
    + \frac{\eta}{\eta + \alpha} \log \pdf_{j}(\vVarCem)
    + C \enspace,
\end{align}
where $C = - \log Z_{\eta, \alpha}$ 
% = -(\nu + \eta + \alpha)/(\eta + \alpha)$ 
absorbs all additive constants. Exponentiating yields the optimal proposal $\pdf_{j+1}^{\ast}$ \cref{eq:opt_proposal}, where the multiplier $\nu$ is implicitly accounted for by $Z_{\eta, \alpha}$ required to satisfy the normalization constraint.


\subsection{Derivation of the Dual Function}
\label[appendix]{sec:appendix:dual}

To derive the dual function $g(\eta, \alpha)$, we substitute the optimal proposal $\pdf_{j+1}^{\ast}$ back into the Lagrangian. 
First, we take the logarithm of the optimal proposal \cref{eq:log_opt_proposal} multiplied by $\eta + \alpha$ (which is strictly positive, by \cref{thm:sum_positive}) to match the coefficient of the entropy term in the Lagrangian.
% \begin{align}
%     \log \pdf_{j+1}^{\ast}(\vVarCem) 
%     = \frac{\eta}{\eta+\alpha} \log \pdf_j(\vVarCem) - \frac{J(\vVarCem)}{\eta+\alpha} - \log Z_{\eta,\alpha} \enspace.
% \end{align}
% Next, we multiply this expression by $(\eta + \alpha)$ to match the coefficient of the entropy term in the Lagrangian
% \begin{align}
%     (\eta + \alpha) \log \pdf_{j+1}^{\ast}(\vVarCem) 
%     & = 
%     \eta \log \pdf_j(\vVarCem) - J(\vVarCem) 
%     \\
%     & \quad
%     - (\eta + \alpha) \log Z_{\eta,\alpha} \enspace.
% \end{align}
Taking the expectation of $(\eta + \alpha) \log \pdf_{j+1}^{\ast}(\vVarCem)$ w.r.t. $\pdf_{j+1}^{\ast}$ yields
\begin{align}
    &(\eta + \alpha) \int_{\IR^{\dimVarCem}} \pdf_{j+1}^{\ast}(\vVarCem) \log \pdf_{j+1}^{\ast}(\vVarCem) \dd \vVarCem 
    = 
    \\
    & \quad
    \eta \int \pdf_{j+1}^{\ast}(\vVarCem) \log \pdf_j(\vVarCem) \dd \vVarCem 
    % &\quad 
    - \int_{\IR^{\dimVarCem}} \pdf_{j+1}^{\ast}(\vVarCem) J(\vVarCem) \dd \vVarCem \\
    &\quad - (\eta + \alpha) \log Z_{\eta,\alpha} \int \pdf_{j+1}^{\ast}(\vVarCem) \dd \vVarCem \enspace.
    \label{eq:appendix:expected_log}
\end{align}
Since $\pdf_{j+1}^{\ast}$ is a valid \ac{PDF}, it integrates to $1$.
% ($\int_{\IR^{\dimVarCem}} \pdf_{j+1}^{\ast} \dd \vVarCem = 1$). 
Then, substituting \cref{eq:appendix:expected_log} back into the simplified Lagrangian reads
\begin{align}
    &\mathcal{L}(\pdf_{j+1}^{\ast},\eta,\alpha) 
    = \int_{\IR^{\dimVarCem}} \pdf_{j+1}^{\ast}(\vVarCem) J(\vVarCem) \dd \vVarCem \nonumber \\
    & \quad + \eta \int_{\IR^{\dimVarCem}} \pdf_{j+1}^{\ast}(\vVarCem) \log \pdf_j(\vVarCem) \dd \vVarCem 
    \\
    & \quad - \int_{\IR^{\dimVarCem}} \pdf_{j+1}^{\ast}(\vVarCem) J(\vVarCem) \dd \vVarCem - (\eta + \alpha) \log Z_{\eta,\alpha} \nonumber \\
    &\quad - \eta \int_{\IR^{\dimVarCem}} \pdf_{j+1}^{\ast}(\vVarCem) \log \pdf_j(\vVarCem) \dd \vVarCem \nonumber \\
    &\quad - \eta \epsilon + \alpha H_{\min} \enspace,
\end{align}
where the normalization multiplier $\nu$ has been absorbed in $Z_{\eta, \alpha}$.
Notice that the terms involving the expected cost 
% $\int_{\IR^{\dimVarCem}} \pdf_{j+1}^{\ast} J \dd \vVarCem$ 
and the cross-entropy $\eta \int_{\IR^{\dimVarCem}} \pdf_{j+1}^{\ast}(\vVarCem) \log \pdf_j(\vVarCem) \dd \vVarCem$ cancel out. 
The remaining terms yield the dual function \cref{eq:tr_dual}.
% \begin{align}
%     g(\eta, \alpha) = -\eta \epsilon + \alpha H_{\min} - (\eta + \alpha) \log Z_{\eta,\alpha} \enspace.
% \end{align}


% \subsection{Gradient of the Dual Function}
% Substituting $\pdf_{j+1}^{\ast}$ from \cref{eq:opt_proposal} into the gradient of the dual function \cref{eq:tr_dual:gradient} gives
% \begin{align}
% & \frac{\partial g}{\partial \eta}
%     =-\log Z_{\eta,\alpha} -\epsilon
%     \\ 
%     &
%     -
%     % \frac{1}{\eta+\alpha}
%     \frac{\Eop_{\pdf_{j}} \left\{\pdf_{j}(\vVarCem)^{\frac{\eta}{\eta+\alpha}-1}\exp \left(-\frac{J(\vVarCem)}{\eta+\alpha}\right)\left(\alpha\log \pdf_{j}(\vVarCem)+J(\vVarCem)\right)\right\}}
%     {(\eta+\alpha) Z_{\eta,\alpha}}
% \\
% & \frac{\partial g}{\partial \alpha}
%     = -\log Z_{\eta,\alpha} +H_{\min}
%     \\
%     & 
%     -
%     % \frac{1}{\eta+\alpha}
%     \frac{\Eop_{\pdf_{j}} \left\{\pdf_{j}(\vVarCem)^{\frac{\eta}{\eta+\alpha}-1}\exp \left(-\frac{J(\vVarCem)}{\eta+\alpha}\right)\left(-\eta\log \pdf_{j}(\vVarCem)+J(\vVarCem)\right)\right\} }
%     {(\eta+\alpha) Z_{\eta,\alpha}}
% \end{align}

% \subsection{Strict Positivity of the Trust Region Multiplier}
% \label[appendix]{sec:appendix:eta_positive}

% \begin{theorem}
%     \label{thm:eta_positive}
%     Let $J$ be non-constant and, if the entropy constraint is active, let $\pdf^{\mathrm{G}} \propto \exp(-\nicefrac{J}{\bar{\alpha}})$ denote the Gibbs \ac{PDF} at the temperature $\bar{\alpha}$ for which the entropy bound is attained, \ie, $H(\pdf^{\mathrm{G}}) = H_{\min}$, and assume $\epsilon < D_{\mathrm{KL}}(\pdf^{\mathrm{G}} \| \pdf_j)$.
%     Then, the trust region is active, \ie, $D_{\mathrm{KL}}(\pdf_{j+1}^{\ast} \| \pdf_j) = \epsilon$, and its multiplier is strictly positive, $\eta^{\ast} > 0$.
% \end{theorem}
% \textit{Proof:} Complementary slackness admits only the two alternatives $\eta^{\ast} = 0$ or $D_{\mathrm{KL}}(\pdf_{j+1}^{\ast} \| \pdf_j) = \epsilon$.
% Assume $\eta^{\ast} = 0$. Then the \ac{KL} term drops out of the Lagrangian~\cref{eq:tr:lagrangian}, whose point-wise stationarity condition $J(\vVarCem) + \nu + \alpha ( \log \pdf_{j+1}(\vVarCem) + 1 ) = 0$ becomes independent of $\pdf_j$.
% For $\alpha^{\ast} = 0$, it reduces to $J(\vVarCem) = -\nu$, contradicting the non-constant $J$; the Lagrangian is then linear in $\pdf_{j+1}$ and minimized by the Dirac delta at $\argmin_{\vVarCem} J(\vVarCem)$, which is not absolutely continuous w.r.t. $\pdf_j$ and thus yields $D_{\mathrm{KL}} = \infty > \epsilon$.
% For $\alpha^{\ast} > 0$, it is solved by $\pdf_{j+1}^{\ast} = \pdf^{\mathrm{G}}$, with $\alpha^{\ast} = \bar{\alpha}$ fixed by the active entropy constraint, so that $D_{\mathrm{KL}}(\pdf^{\mathrm{G}} \| \pdf_j) > \epsilon$ by assumption.
% In both cases, the candidate is infeasible, contradicting optimality.
% Hence, only the second alternative remains, which requires $\eta^{\ast} > 0$. \hfill$\blacksquare$

% Since $H(\pdf^{\mathrm{G}})$ increases strictly with the temperature, $\bar{\alpha}$ is unique.
% The assumption on $\epsilon$ can be certified in advance.
% Splitting the \ac{KL} divergence into entropy and cross-entropy and using $H(\pdf^{\mathrm{G}}) = H_{\min}$ gives $D_{\mathrm{KL}}(\pdf^{\mathrm{G}} \| \pdf_j) = -H_{\min} - \Eop_{\pdf^{\mathrm{G}}}\{ \log \pdf_j(\vVarCem) \}$.
% Bounding the log-density by its maximum, $\log \pdf_j(\vVarCem) \leq \log \sup_{\vVarCem} \pdf_j(\vVarCem)$, then yields $D_{\mathrm{KL}}(\pdf^{\mathrm{G}} \| \pdf_j) \geq -H_{\min} - \log \sup_{\vVarCem} \pdf_j(\vVarCem)$.
% For a Gaussian proposal, $\sup_{\vVarCem} \pdf_j(\vVarCem) = ((2\pi)^{\dimVarCem} \det \mC_j)^{-\nicefrac{1}{2}}$, so that this lower bound becomes $H(\pdf_j) - H_{\min} - \nicefrac{\dimVarCem}{2}$, \ie, it suffices to choose the entropy floor sufficiently far below the entropy of the current proposal.
% Conversely, if the assumption is violated, $\pdf^{\mathrm{G}}$ lies within the trust region and is thus already optimal, \ie, $\alpha^{\ast} = \bar{\alpha} > 0$ and $\eta^{\ast} = 0$, which shows that the assumption can only fail if the entropy constraint is present.
% Strict positivity of $\eta^{\ast}$ implies that the exponent $\nicefrac{\eta^{\ast}}{(\eta^{\ast}+\alpha^{\ast})}$ in \cref{eq:opt_proposal} is non-zero, \ie, the previous proposal is always retained.

% Independently of this assumption, the update \cref{eq:opt_proposal} is always well-defined.

\subsection{Strict Positivity of the Lagrange Multiplier Sum}
\label[appendix]{sec:appendix:pos_multiplier_sum}
\begin{theorem}
    \label{thm:sum_positive}
    Let $J$ be non-constant, $H_{\min} > -\infty$, $f_j$ be non-degenerate, and $\epsilon < \infty$. Then $\eta^{\ast} + \alpha^{\ast} > 0$.
\end{theorem}
% \textit{Proof:} 
\begin{proof}
By \ac{KKT} dual feasibility, $\eta^{\ast} \geq 0$ and $\alpha^{\ast} \geq 0$, so $\eta^{\ast} + \alpha^{\ast} = 0$ requires $\eta^{\ast} = \alpha^{\ast} = 0$.
Under this assumption, the Lagrangian~\cref{eq:tr:lagrangian} reduces to the linear functional $\int_{\IR^{\dimVarCem}} \pdf_{j+1}(\vVarCem) J(\vVarCem) \dd \vVarCem + \nu \left( \int_{\IR^{\dimVarCem}} \pdf_{j+1}(\vVarCem) \dd \vVarCem - 1 \right)$,
and \cref{eq:lagrange_func_derivative} reduces to $\nicefrac{\delta \mathcal{L}}{\delta \pdf_{j+1}} = J(\vVarCem) + \nu$.
Solving for $\pdf_{j+1}$ over valid probability measures requires either:
% \begin{enumerate}[(i)]
i) $J(\vVarCem) = -\nu$ everywhere on the support of $\pdf_{j+1}$, which requires $J$ to be constant if $\pdf_{j+1}$ is continuous and strictly positive; or
ii) $\pdf_{j+1} = \delta(\vVarCem - \vVarCem^*)$ with $\vVarCem^* \in \argmin_{\vVarCem} J(\vVarCem)$, concentrating all mass onto the global minimum of $J$.
% \end{enumerate}
The first case directly contradicts the assumption that $J$ is non-constant. The second case yields $H(\delta) = -\infty < H_{\min}$ and $D_{\mathrm{KL}}(\delta \| \pdf_j) = \infty > \epsilon$, where either violation suffices to contradict primal feasibility and \ac{KKT} optimality.
Therefore, $\eta^{\ast} = \alpha^{\ast} = 0$ cannot hold, implying $\eta^{\ast} + \alpha^{\ast} > 0$. %\hfill$\blacksquare$
\end{proof}
Hence, the exponents in \cref{eq:opt_proposal} never degenerate, regardless of which of the two constraints is active.


\subsection{Strong Duality}
\label[appendix]{sec:appendix:strong_duality}
% argumentations, 
% 1. expexted cost is linear in $\pdf_{j+1}$,teherfore the primal problem is a convex 
% 2. the equality constraint is linear in $\pdf_{j+1}$,
% 3. inequalities (reformulated in standard form) $D_{\mathrm{KL}}(\pdf_{j+1} \| \pdf_j) - \epsilon \leq 0$ and $H_{\min} - H(\pdf_{j+1}) \leq 0$ are strictly convex functionals of $\pdf_{j+1}$, since $D_{\mathrm{KL}}(\pdf_{j+1} \| \pdf_j)$ and $-H(\pdf_{j+1})$ are strictly convex functionals of $\pdf_{j+1}$. subtracting the constant $\epsilon$ and $H_{\min}$ does not change the convexity of the functionals.
% 4. there exists a strictly feasible point, namely $\pdf_{j+1} = \pdf_j$, since $D_{\mathrm{KL}}(\pdf_j \| \pdf_j) = 0 < \epsilon$ and $H(\pdf_j) > H_{\min}$, the latter holding by construction since the entropy floor is chosen below the entropy of the current proposal.
% therefore Slater's condition is satisfied, and strong duality holds, i.e., minimizing the dual function recovers the exact primal optimum.

% Slater's condition~\cite[Sec.~5.2.3]{boydConvex2004} guarantees strong duality for a problem whose objective function and inequality constraints are convex and whose equality constraints are linear, provided that there exists a strictly feasible point, that is, a point at which all inequality constraints hold strictly (excluding equality).
% It thus remains to verify these three properties.

% Both the expected cost and the normalization constraint are linear in $\pdf_{j+1}$ in \cref{eq:tr_objective_function}, therefore, the objective function is convex with a linear equality constraint.
% %
% Written in standard form, the inequality constraints $D_{\mathrm{KL}}(\pdf_{j+1} \| \pdf_j) - \epsilon \leq 0$ and $H_{\min} - H(\pdf_{j+1}) \leq 0$ are strictly convex functions of $\pdf_{j+1}$. %, as subtracting the constants $\epsilon$ and $H_{\min}$ does not affect convexity.
% The primal problem is therefore a convex program. % over a convex feasible set.
% Moreover, the point $\pdf_{j+1} = \pdf_j$ is strictly feasible, since $D_{\mathrm{KL}}(\pdf_j \| \pdf_j) = 0 < \epsilon$ and $H(\pdf_j) > H_{\min}$, the latter holding by construction as the entropy floor is chosen below the entropy of the current proposal.
% Slater's condition is thus satisfied, and strong duality holds. %, \ie, the duality gap vanishes, $g(\eta^{\ast}, \alpha^{\ast}) = \Eop_{\pdf_{j+1}^{\ast}}\{J(\vVarCem)\}$.
% Consequently, solving the dual and evaluating \cref{eq:opt_proposal} at $(\eta^{\ast}, \alpha^{\ast})$ recovers the exact primal optimum.
%
%
We show strong duality under the following two weak assumptions, which hold for common setups such as positive-definite quadratic costs and Gaussian proposals.

\begin{assumption}
    \label{assumption:duality:on_cost}
    We consider minimizing the expected cost over a class of probability distributions $\mathcal{F}$.
    We assume that $J(\vVarCem)\ge c > -\infty,~ \forall \vVarCem \in \operatorfont{dom} (\mathcal{F})$ is chosen so that its expectation is bounded below by $c$ across the domain $\operatorfont{dom}(\mathcal{F})$ of the considered distributions in $\mathcal{F}$.
    This holds trivially without restricting $\mathcal{F}$ whenever $J(\vVarCem) \ge c > -\infty,~\forall \vVarCem$.
\end{assumption}
\begin{assumption}
    \label{assumption:duality:on_cost_and_proposal}
    The cost function and initial proposal distribution $\pdf_j \in \mathcal{F}$ are chosen such that $\Eop_{\pdf_j}\{J(\vVarCem)\} < +\infty$ holds.
\end{assumption}
% In practice, both assumptions are weak and hold for common setups such as positive-definite quadratic costs and Gaussian proposals.

To establish strong duality for \cref{eq:tr_objective_function}, we apply Luenberger's duality theorem for infinite-dimensional convex programs~\cite[Sec.~8.6, Thm.~1]{luenbergerOptimizationVectorSpace1969}, which generalizes Slater's condition from finite-dimensional vector spaces to function spaces. Under this theorem, strong duality holds if 1) the cost function and inequality constraints are convex, 2) there exists a strictly feasible point in $\pdf_{j+1}\in\mathcal{F}$, that is, there is a point where all inequality constraints hold strictly (excluding equality), and 3) the primal infimum ($d^*=\inf_{\pdf_{j+1\in \mathcal{F}, G\le0}} \Eop_{\pdf_{j+1}}\{J(\vVarCem)\}$) over the feasible domain of valid probability densities $\mathcal{F}$ is finite, where $G$ summarizes all inequality constraints.

We verify these requirements for \cref{eq:tr_objective_function} over $\mathcal{F}$, and by restricting $\mathcal{F}$ to valid distributions, we implicitly address the equality constraint for normalized \acp{PDF}.
%$\Omega = \{ \pdf \in L^1(\mathcal{X}) \mid \pdf(x) \ge 0 \text{ a.e.}, \, \mathbb{E}_{\pdf}\{1\} = 1 \}$:
\begin{enumerate}
    \item 
    \textbf{Convex structure:} Both the expected cost $\Eop_{\pdf_{j+1}}\{J(\vVarCem)\}$ and the normalization constraint are linear in $\pdf_{j+1}$. The inequality constraints $D_{\mathrm{KL}}(\pdf_{j+1} \Vert{} \pdf_j) - \epsilon \le 0$ and $H_{\min} - H(\pdf_{j+1}) \le 0$ are strictly convex functionals of $\pdf_{j+1}$, yielding convex feasible sets. The primal problem is thus a convex program.
    \item
    \textbf{Strict feasibility:} Setting $\pdf_{j+1} = \pdf_j \in \mathcal{F}$ yields a strictly feasible point, as $D_{\mathrm{KL}}(\pdf_j \Vert{} \pdf_j) = 0 < \epsilon$ and $H(\pdf_j) > H_{\min}$, where the latter holds by construction since the entropy floor is chosen strictly below the entropy of the current proposal.
    \item 
    \textbf{Finite primal infimum:} The primal infimum $d^*$ is finite due to three factors: 
    (i) the feasible set is non-empty (guaranteed by $\pdf_{j+1} = \pdf_j$);
    (ii) $+\infty$ is excluded since $\pdf_j$ satisfies $\Eop_{\pdf_j}\{J(\vVarCem)\} < +\infty$ by \cref{assumption:duality:on_cost_and_proposal}; and
    (iii) the expectation is bounded from below by \cref{assumption:duality:on_cost}.
\end{enumerate}

Since all conditions are satisfied in our setting, Luenberger's theorem for strong duality holds. 
% Consequently, strong duality is established, and optimal non-negative dual multipliers exist.

